\documentclass[11pt,a4paper]{article}
\usepackage[a4paper,top=10mm,bottom=10mm,left=13mm,right=13mm]{geometry}
\usepackage[T1]{fontenc}
\usepackage{amsmath}
\usepackage{amssymb}
\usepackage{enumitem}
\usepackage{multicol}
\pagestyle{empty}
\setlength{\parindent}{0pt}
\setlist[enumerate,1]{label=\textbf{Q\arabic*.}, leftmargin=*, itemsep=5pt, topsep=3pt, parsep=0pt}
\setlist[enumerate,2]{label=(\alph*), leftmargin=*, itemsep=2pt, topsep=1pt, parsep=0pt}

\begin{document}

\begin{center}
  {\Large\bfseries Percentages: Non-calculator questions --- Answer Key}
\end{center}

\begin{enumerate}

\item
\begin{enumerate}
  \item \(51\%\)
  \item \(0.51\)
  \item \(\frac{51}{100}\)
\end{enumerate}

\item
\begin{enumerate}
  \item \(0.4=40\%\)
  \item \(0.06=6\%\)
  \item \(1.25=125\%\)
  \item \(45\%=0.45\)
  \item \(7\%=0.07\)
  \item \(130\%=1.3\)
\end{enumerate}

\item
\begin{enumerate}
  \item Neon: \(36\%\)
  \item Jungle: \(10\%\)
  \item \(54\%\); yes, Ravi is right.
\end{enumerate}

\item
\begin{enumerate}
  \item \(25\)
  \item \(3\)
  \item \(28\)
\end{enumerate}

\item
\begin{enumerate}
  \item Sale price: \pounds36
  \item New price: \pounds40
\end{enumerate}

\item
\begin{enumerate}
  \item Original price: \pounds60
  \item Membership before the rise: \pounds70 per month
\end{enumerate}

\item VAT: \pounds11.20

\item Cold drinks sold: \(108\)

\item No. A \(20\%\) increase followed by a \(20\%\) decrease leaves the price at \(96\%\) of the original; for example, \pounds100 becomes \pounds120 and then \pounds96.

\item
\begin{enumerate}
  \item Shop A: \pounds56;\quad Shop B: \pounds52.\quad Shop B is cheaper by \pounds4.
  \item One possible third offer: \(20\%\) off, then \(\pounds15\) off. This gives \pounds49, which is cheaper than both shops and is more than half price.
  \item \pounds160\ (both shops would charge \pounds112)
\end{enumerate}

\end{enumerate}

\end{document}